\chapter*{三、研究方法}

\section*{3.1 本章架構與方法總覽}

本章的內容安排可分為三個部分。第3.2節至第3.6節說明本文提出之
OOT-AEnbPI及其Local-scale延伸：第3.2節先介紹整體流程，第3.3節與
第3.4節說明Kalman高低頻分解及基礎點預測器，第3.5節說明
OOT-AEnbPI原始區間，第3.6節說明Local-scale延伸。第3.7節以
IFOMC作為替代區間預測器；第3.8節與第3.9節介紹B1至B5五種比較
方法及其Local-scale校準方式。第3.10節定義本文使用的評估指標。

表\ref{tab:method-overview}整理本文提出方法、延伸區間版本與比較
方法的整體架構。OOT-AEnbPI、OOT-AEnbPI-LS及
OOT-AEnbPI-IFOMC使用相同的Kalman--ARIMA--ANN ensemble基礎點
預測，三者的差異集中於區間預測器。B1至B5則沿用既有研究的高頻
模型與區間生成方式，作為評估OOT-AEnbPI的比較基準。

\begin{table}[htbp]
\centering
\caption{本文提出方法與比較方法之整體架構}
\label{tab:method-overview}
\resizebox{\textwidth}{!}{%
\begin{tabular}{clll}
\toprule
定位 & 方法 & 點預測架構 & 區間預測方式\\
\midrule
本文核心方法 & OOT-AEnbPI
& Kalman--ARIMA--ANN ensemble
& OOT residual、自適應window與最短非對稱區間\\
本文延伸版本 & OOT-AEnbPI-LS
& Kalman--ARIMA--ANN ensemble
& 局部尺度標準化後的OOT residual區間\\
替代區間比較 & OOT-AEnbPI-IFOMC
& Kalman--ARIMA--ANN ensemble
& OOT residual狀態轉移區間\\
\midrule
比較方法B1 & DistPred--ANN
& Kalman--ARIMA--ANN
& DistPred經驗分位數區間\\
比較方法B2 & DistPred--Minusformer
& Kalman--ARIMA--Minusformer
& DistPred經驗分位數區間\\
比較方法B3 & Quantile--ANN
& Kalman--ARIMA--ANN
& Quantile條件分位數區間\\
比較方法B4 & Quantile--Minusformer
& Kalman--ARIMA--Minusformer
& Quantile條件分位數區間\\
比較方法B5 & Sort Bootstrap--ANN
& Kalman--ARIMA--ANN
& Bootstrap預測的經驗分位數區間\\
\bottomrule
\end{tabular}%
}
\end{table}

\section*{3.2 演算法流程}

本研究提出Out-of-Time Adaptive EnbPI（以下簡稱OOT-AEnbPI）。OOT-AEnbPI以
Kalman--ARIMA--ANN ensemble作為基礎點預測器，並針對時間序列
預測區間所使用的校準殘差進行時間順序調整。利用依時間順序產生的out-of-time（OOT）residual
建立殘差池，只使用training期間的資訊選擇residual window，再由
選定的近期殘差建立可逐期更新的最短非對稱預測區間。

OOT-AEnbPI的主要流程如下：
\begin{enumerate}
  \item 使用Kalman filter將觀察序列分解為低頻與高頻成分；
  \item 分別以ARIMA及ANN ensemble預測低頻與高頻成分，再合併為原序列的點預測；
  \item 將training前半段設為initial prefix，其餘資料依時間順序劃分為多個future blocks；
  \item 對每個future block只使用該block以前的資料估計模型，再預測該block，以建立OOT residual pool；
  \item 在training期間以逐期回測比較候選residual windows；
  \item 從選定的殘差池計算最短非對稱上下界，建立OOT-AEnbPI原始區間；
  \item 發出當期預測後才觀察真值，並更新殘差池供下一期使用。
\end{enumerate}

在上述OOT-AEnbPI核心方法之外，本研究另建立兩種區間版本。第一種為
OOT-AEnbPI-LS，將OOT residual依當期局部狀態進行尺度標準化；
第二種為OOT-AEnbPI-IFOMC，利用前一期殘差狀態及逐期更新的
轉移矩陣建立條件區間。兩者均保留OOT-AEnbPI的基礎點預測，其中IFOMC主要
用於比較不同區間預測器，不屬於OOT-AEnbPI原始區間的核心流程。

為評估OOT-AEnbPI，本研究另與\cite{hung2026thesis}及\cite{lin2024thesis}所提出的方法進行比較。
將比較方法依序記為B1至B5。B1至B5沿用既有模型定義及原始超參數，其實際設定統一於本章後段說明。

\section*{3.3 Kalman分解}

先使用一維local-level Kalman filter處理觀察序列$\{z_t\}$。令$x_t$為無法直接觀察的潛在水準狀態，
模型的狀態方程與觀測方程分別為
\begin{align}
  x_t &= x_{t-1}+w_{t-1}, & w_{t-1}&\sim N(0,Q),\\
  z_t &= x_t+v_t, & v_t&\sim N(0,R),
\end{align}
其中$w_t$為process noise、$v_t$為measurement noise，並假設兩者互相獨立。
$Q$控制潛在水準允許隨時間改變的程度，$R$則表示觀察值相對於潛在水準的雜訊程度。

使用第一筆觀察值及初始估計誤差變異數初始化：
\begin{equation}
 \widehat x_{0|0}=z_0,\qquad P_{0|0}=1.
\end{equation}
接著依時間順序逐期執行預測與更新。首先，在尚未使用第$t$期觀察值以前，由前一期更新後的估計得到當期先驗狀態：
\begin{align}
 \widehat x_{t|t-1}&=\widehat x_{t-1|t-1},\\
 P_{t|t-1}&=P_{t-1|t-1}+Q.
\end{align}

取得第$t$期觀察值$z_t$後，計算Kalman gain：
\begin{equation}
 K_t=\frac{P_{t|t-1}}{P_{t|t-1}+R}.
\end{equation}
再利用當期innovation $z_t-\widehat x_{t|t-1}$更新水準估計與其誤差變異數：
\begin{align}
 \widehat x_{t|t}
 &=\widehat x_{t|t-1}
 +K_t\left(z_t-\widehat x_{t|t-1}\right),\\
 P_{t|t}&=(1-K_t)P_{t|t-1}.
\end{align}

完成第 $t$ 期的更新後，以 Kalman filter 得到的狀態估計值
$\widehat{x}_{t|t}$ 作為變動較平緩的水準成分，並將觀察值與
狀態估計值之差作為變動較快速的剩餘成分：
\begin{align}
    \ell_t &= \widehat{x}_{t|t},\\
    r_t &= z_t-\widehat{x}_{t|t}.
\end{align}
令水準成分與剩餘成分在第$t-1$期資訊下，對第$t$期的
一步點預測分別為$\widehat{\ell}_{t|t-1}$與
$\widehat{r}_{t|t-1}$，則原觀察序列的一步合併點預測為
\begin{equation}
    \widehat{z}_{t|t-1}
    =
    \widehat{\ell}_{t|t-1}
    +
    \widehat{r}_{t|t-1}.
\end{equation}
取得合併點預測後，再依不同區間估計方法計算相對於
$\widehat{z}_{t|t-1}$ 的下界與上界偏移量。令兩個偏移量分別為
$q_t^{-}$ 與 $q_t^{+}$，則最終預測區間可統一表示為
\begin{equation}
    C_t
    =
    \left[
    \widehat{z}_{t|t-1}+q_t^{-},
    \widehat{z}_{t|t-1}+q_t^{+}
    \right].
\end{equation}


\section*{3.4 OOT-AEnbPI基礎點預測器}

\subsection*{3.4.1 低頻ARIMA}

令Kalman filter得到的低頻成分為
$\{\ell_1,\ldots,\ell_n\}$。OOT-AEnbPI使用$ARIMA(p,0,q)$描述低頻成分的
線性時間相依性。由於低頻成分已是變動較平緩的水準序列，差分
階數固定為$d=0$。自迴歸階數與移動平均階數的候選集合為
\[
p,q\in\{0,1,2,3,4\},\qquad (p,q)\neq(0,0).
\]
對每一組候選$(p,q)$估計$ARIMA(p,0,q)$，再以BIC選擇階數：
\[
(\widehat p,\widehat q)
=
\arg\min_{p,q}
BIC\{ARIMA(p,0,q)\}.
\]

選定階數後，以當時所有可用的低頻序列重新估計模型。預測第$t$期
時，ARIMA只能使用第$t-1$期以前已取得的低頻成分，產生一步向前
點預測
\[
\widehat\ell_{t|t-1}
=
E(\ell_t\mid\ell_1,\ldots,\ell_{t-1}).
\]
在建立training期future-block OOT residual時，每個block的ARIMA
階數選擇與參數估計皆只使用該block以前的資料；在testing開始前，
則以完整training低頻序列重新選階與估計。如此可避免使用future
block或testing期間的未來資訊。

\subsection*{3.4.2 高頻ANN}

OOT-AEnbPI的ANN使用長度$w=15$的滑動視窗。預測第$t$期高頻成分$r_t$
時，輸入同時包含第$t$期以前的低頻成分、高頻成分及原觀察序列：
\begin{equation}
 X_t=\left(
 \ell_{t-15:t-1},\;r_{t-15:t-1},\;z_{t-15:t-1}
 \right),
\end{equation}
因此每筆輸入共有45個特徵，預測目標為$r_t$。所有輸入都只到
$t-1$期為止，不包含當期或未來資料。

ANN包含兩個隱藏層$(32,16)$，使用ReLU及Adam；learning rate為
0.001、正則化係數為$10^{-4}$。輸入特徵與高頻目標值均使用當次
training資料估計的平均數與標準差進行標準化。最大迭代次數的候選
集合為$\{125,250,500\}$，並以3個時間順序rolling validation
splits的平均MSE選擇；每個split只使用較早資料訓練、較晚資料驗證。
選定後，再以當時完整可用的training樣本重新估計ANN。

\subsection*{3.4.3 Bootstrap hybrid ensemble}

為降低單一模型的不穩定性，OOT-AEnbPI建立$B=30$組bootstrap hybrid
models。設訓練樣本數為$n$，區塊長度為
\begin{equation}
 d_B=\max\left\{2,\operatorname{round}(\sqrt n)\right\}.
\end{equation}
第$b$次bootstrap隨機抽取區塊起點$S_{b,j}$，並以
\begin{equation}
 \mathcal I_{b,j}=\left\{(S_{b,j}+k)\bmod n:
 k=0,\ldots,d_B-1\right\}
\end{equation}
取得連續樣本索引，直到組成$n$筆樣本，因此屬於circular
moving-block bootstrap。

低頻部分保留原始時間路徑，只依上述索引重抽中心化後的一步預測
殘差，避免拼接低頻區塊造成不自然的水準跳動；高頻部分則使用相同
索引重抽輸入與目標。每組樣本分別估計ARIMA與ANN，形成一個
hybrid model。令第$b$組模型的高低頻預測為
$\widehat\ell_{t|t-1}^{(b)}$與$\widehat r_{t|t-1}^{(b)}$，則原始
ensemble點預測為
\begin{equation}
 \widehat z_{t|t-1}^{raw}
 =\frac{1}{B}\sum_{b=1}^{B}
 \left(\widehat\ell_{t|t-1}^{(b)}+
 \widehat r_{t|t-1}^{(b)}\right).
\end{equation}

\section*{3.5 OOT-AEnbPI原始區間}

\subsection*{3.5.1 OOT residual的建立}

OOT-AEnbPI將training前$\rho=0.5$比例設為initial prefix，剩餘資料按時間
順序分為4個future blocks。對第$k$個block，模型只使用該block
以前的資料估計，再預測該block，因此得到
\begin{equation}
 e_t^{OOT}=z_t-\widetilde z_{t|t-1},
\end{equation}
其中$\widetilde z_{t|t-1}$表示未使用第$t$期及其未來資料所產生的
預測。依序合併各block的殘差形成OOT residual pool。這種設計目的是使校準殘差更接近真正測試時
可取得的樣本外誤差。

\subsection*{3.5.2 自適應residual window}

OOT-AEnbPI原始區間的候選集合為
\begin{equation}
 \mathcal W=\{100,200,300,400,\mathrm{all}\}.
\end{equation}
每個候選$W$只保留最近$W$筆OOT residual，並在OOT residual序列
尾端40\%進行prequential回測。預測某一驗證殘差時，只能使用它
以前的殘差建立區間；以平均interval score最小的候選值作為
$\widehat W$。所有候選比較及視窗選擇均限制於training期間完成。

\subsection*{3.5.3 Bias correction與最短非對稱區間}

令第$t$期可用的raw residual pool為$\mathcal E_t$。OOT-AEnbPI使用其平均值
進行bias correction：
\begin{equation}
 b_t=\frac{1}{|\mathcal E_t|}\sum_{e_i\in\mathcal E_t}e_i,
 \qquad
 \widehat z_{t|t-1}^{\mathrm{OOT-AEnbPI}}=\widehat z_{t|t-1}^{raw}+b_t.
\end{equation}
為避免相同偏差同時出現在點預測與區間偏移量，分位數計算使用
中心化殘差$\mathcal E_t-b_t$。

OOT-AEnbPI不強制左右各保留$\alpha/2$的誤差，而是在
$\beta\in[0,\alpha]$設定101個等距網格點，尋找
\begin{equation}
 \widehat\beta_t=\arg\min_{\beta}
 \left\{Q_{1-\alpha+\beta}(\mathcal E_t-b_t)
       -Q_{\beta}(\mathcal E_t-b_t)\right\}.
\end{equation}
令$q_t^-=Q_{\widehat\beta_t}(\mathcal E_t-b_t)$、
$q_t^+=Q_{1-\alpha+\widehat\beta_t}(\mathcal E_t-b_t)$，則
\begin{equation}
 C_t^{\mathrm{OOT-AEnbPI}}=\left[
 \widehat z_{t|t-1}^{\mathrm{OOT-AEnbPI}}+q_t^-,
 \widehat z_{t|t-1}^{\mathrm{OOT-AEnbPI}}+q_t^+
 \right].
\end{equation}
完成第$t$期點預測與區間後才觀察$z_t$，再將raw residual
$z_t-\widehat z_{t|t-1}^{raw}$加入pool，移除超出$\widehat W$的
舊殘差，供第$t+1$期使用。

\section*{3.6 OOT-AEnbPI-LS}

OOT-AEnbPI-LS保留OOT-AEnbPI的OOT residual建立方式，但讓殘差尺度依局部
資料狀態調整，以避免高波動與低波動狀態共用完全相同的誤差尺度。

\subsection*{3.6.1 狀態特徵與局部尺度}

對時間點$i$建立四維特徵
\begin{equation}
 X_i=\left(
 \widehat z_{i|i-1},
 \widehat z_{i|i-1}-z_{i-1},
 SD(\Delta z_{i-12:i-1}),
 SD(\Delta z_{i-36:i-1})
 \right).
\end{equation}
四項依序為預測水準、相對前一期的變化、最近12期一階差分的短期
波動，以及最近36期一階差分的長期波動；12與36是波動觀察窗，
不是季節差分。

各特徵先依training資料的中位數與四分位距作RobustScaler轉換，
再以標準化特徵的平方Euclidean distance尋找最近的$K=60$個歷史
時間點。令鄰居集合為$N_{60}(X_t)$，局部尺度定義為
\begin{equation}
 \widehat\sigma_t=\operatorname{median}
 \left\{|z_i-\widehat z_{i|i-1}|:i\in N_{60}(X_t)\right\},
\end{equation}
並使用
\begin{equation}
 \widehat\sigma_t\leftarrow
 \max\{\widehat\sigma_t,Q_{0.10}(|e|),10^{-6}\}
\end{equation}
避免尺度過小。training期估計尺度時排除樣本本身；testing期只能從
當時已觀察的歷史點尋找鄰居。當期真值出現後，才將其特徵與絕對
殘差加入資料庫。

\subsection*{3.6.2 標準化殘差與區間生成}

對每筆OOT residual計算
\begin{equation}
 u_i=\frac{e_i^{OOT}}{\widehat\sigma_i}.
\end{equation}
Local-scale版本在標準化殘差上，從
\begin{equation}
 \mathcal W_{LS}=\{24,60,120,180,300,\mathrm{all}\}
\end{equation}
選擇window。令當期標準化殘差pool為$\mathcal U_t$，其平均為
\begin{equation}
 b_t^u=\frac{1}{|\mathcal U_t|}\sum_{u_i\in\mathcal U_t}u_i.
\end{equation}
修正後點預測為
\begin{equation}
 \widehat z_{t|t-1}^{LS}
 =\widehat z_{t|t-1}^{raw}+\widehat\sigma_t b_t^u.
\end{equation}
對$u_i-b_t^u$尋找95\%最短非對稱分位數端點
$q_t^{u,-}$與$q_t^{u,+}$，回到原尺度後得到
\begin{equation}
 C_t^{\mathrm{OOT-AEnbPI-LS}}=\left[
 \widehat z_{t|t-1}^{LS}+\widehat\sigma_tq_t^{u,-},
 \widehat z_{t|t-1}^{LS}+\widehat\sigma_tq_t^{u,+}
 \right].
\end{equation}
觀察$z_t$後，以raw point計算
\begin{equation}
 u_t^{new}=\frac{z_t-\widehat z_{t|t-1}^{raw}}{\widehat\sigma_t},
\end{equation}
再加入標準化residual pool，供下一期使用。

\section*{3.7 OOT-AEnbPI-IFOMC}

為比較不同區間預測器，本研究保留OOT-AEnbPI已完成bias correction的點
預測，另將\cite{dingmeng2020}的IFOMC狀態轉移概念應用
於OOT-AEnbPI的OOT residual。此版本不使用原文的EMD remaining component，
而是將OOT-AEnbPI依時間順序建立的誤差視為Markov狀態序列。

\subsection*{3.7.1 殘差狀態與轉移矩陣}

令raw OOT residual為$e_i^{OOT}$。為與OOT-AEnbPI testing點預測的bias
correction一致，先扣除當時只由過去residual pool計算的平均值：
\begin{equation}
 r_i=e_i^{OOT}-\overline e_{i-1}.
\end{equation}
將training期$r_i$由小到大以經驗分位數切成$K=10$個樣本數盡量
接近的有序狀態$S_1,\ldots,S_K$。狀態邊界只使用training OOT
residual估計，testing期間固定不變。令$m_{ab}$為相鄰時間由
$S_a$轉移至$S_b$的次數，則
\begin{equation}
 p_{ab}=\frac{m_{ab}}{\sum_{j=1}^{K}m_{aj}}.
\end{equation}

\subsection*{3.7.2 條件區間與rolling更新}

預測第$t$期時只能使用前一期已觀察殘差狀態$s_{t-1}=S_a$。
轉移矩陣第$a$列描述$r_t$進入各目的狀態的條件機率；實作將
$p_{aj}$平均分配給歷史上屬於$S_j$的殘差值，形成加權條件殘差
分布，並取
\begin{equation}
 q_t^-=Q_{0.025}(r_t\mid s_{t-1}=S_a),\qquad
 q_t^+=Q_{0.975}(r_t\mid s_{t-1}=S_a).
\end{equation}
若第$a$列累積轉移少於10次，則改用當時所有可得的歷史殘差，避免
以極少轉移估計尾端分位數。IFOMC區間為
\begin{equation}
 C_t^{\mathrm{OOT-AEnbPI-IFOMC}}=\left[
 \widehat z_{t|t-1}^{\mathrm{OOT-AEnbPI}}+q_t^-,
 \widehat z_{t|t-1}^{\mathrm{OOT-AEnbPI}}+q_t^+
 \right].
\end{equation}
發出區間後才觀察$z_t$並計算
\begin{equation}
 r_t=z_t-\widehat z_{t|t-1}^{\mathrm{OOT-AEnbPI}}.
\end{equation}
再依固定邊界取得$s_t$，新增$s_{t-1}\rightarrow s_t$轉移，供下一期
使用。IFOMC不屬於conformal prediction，本文僅報告testing period
的實證coverage與interval score，不主張嚴格有限樣本保證。

\section*{3.8 比較方法B1至B5}

B1至B5是用來評估OOT-AEnbPI的比較基準。五種方法皆先使用第3.3節的
Kalman filter分解資料，低頻成分使用$ARIMA(p,0,q)$；候選
$p,q\in\{0,1,2,3,4\}$並以AIC選階。各方法的差異集中於高頻點
預測器與區間估計器，其整體組合已整理於表\ref{tab:method-overview}，
以下分別說明各模型及區間方法的實際設定。

\subsection*{3.8.1 人工神經網路（ANN）}

ANN以過去30期高頻值作為輸入，預測下一期高頻成分。網路包含
三個全連接隱藏層，神經元數依序為$(128,256,128)$；每層使用ReLU
activation，並加入dropout 0.2以降低過度配適。訓練使用Adam，
learning rate為0.001、batch size為32、最多訓練50 epochs，early
stopping的patience設為5。

使用DistPred或Quantile輸出時，可用樣本隨機切為80\% training與
20\% validation，並在validation loss連續5個epochs未改善時停止；
使用Sort Bootstrap時，則先對每組重抽樣序列建立滑動視窗，再依
形成後的樣本順序以前80\%訓練、後20\%驗證。ANN隱藏層設定相同，
差異僅在輸出維度及損失函數。

\subsection*{3.8.2 Minusformer}

Minusformer同樣使用過去30期高頻值，label length為15，輸入先經embedding
投影至$d_{model}=128$，包含2個encoder layers、feed-forward
dimension 128、dropout 0.1並啟用gate。每層抽取已解釋訊號$x_l$後
以
\begin{equation}
 h_{l+1}=h_l-x_l
\end{equation}
更新剩餘訊號，使下一層繼續處理尚未被前一層解釋的部分。與一般
Transformer以殘差加法保留輸入不同，Minusformer的input stream
逐層扣除已抽取訊號；output stream則將各層抽取到的資訊轉換為
預測貢獻並逐層累加。因此，減法發生於待解釋的hidden state，並非
直接把前一層的最終預測值從下一層預測值扣除。

每個encoder layer先對輸入進行normalization，再以卷積式
feed-forward模組抽取$x_l$。Gate使用Sigmoid函數控制資訊通過比例，
可概念化表示為
\begin{align}
 g_l&=\sigma\!\left(\operatorname{Conv}_g(h_l)\right),\\
 o_l&=g_l\odot\operatorname{Conv}_o(x_l),
\end{align}
其中$g_l$為門控權重、$o_l$為第$l$層的預測貢獻。各層$o_l$相加
後再經projection映射至需要的輸出維度。搭配DistPred時輸出100個
members；搭配Quantile時輸出三個quantile raw values。兩者使用
相同Minusformer骨幹，差別僅在projection輸出維度及訓練損失。

\subsection*{3.8.3 分布預測（DistPred）}

DistPred不預先假設高頻成分服從常態分布，而是在一次forward pass
輸出$K=100$個DistPred members，以其經驗分布近似下一期高頻成分
的條件分布。這100個輸出可由ANN輸出層或Minusformer最後的
projection產生，並以ensemble CRPS訓練：
\begin{equation}
 \widehat{CRPS}
 =\frac{1}{K}\sum_{k=1}^{K}
  \left|\widehat r^{(k)}-r\right|
 -\frac{1}{2K^2}\sum_{j=1}^{K}\sum_{k=1}^{K}
  \left|\widehat r^{(j)}-\widehat r^{(k)}\right|.
\end{equation}
第一項衡量預測members與真值的距離，使預測分布位於適當位置；
第二項反映members彼此的分散程度，避免100個輸出全部塌縮於相同
數值。程式先將members排序，再使用probability weighted moments
的等價形式降低兩兩比較的計算量。

令排序後輸出為
\begin{equation}
 \widehat r_{t|t-1}^{(1)}\leq\cdots\leq
 \widehat r_{t|t-1}^{(100)},
\end{equation}
高頻點預測取中位數，高頻區間取2.5與97.5百分位數：
\begin{align}
 \widehat r_{t|t-1}&=Q_{0.5}
 \left(\{\widehat r_{t|t-1}^{(k)}\}_{k=1}^{100}\right),\\
 [\widehat L_t^{(H)},\widehat U_t^{(H)}]
 &=\left[Q_{0.025}
 \left(\{\widehat r_{t|t-1}^{(k)}\}_{k=1}^{100}\right),
 Q_{0.975}
 \left(\{\widehat r_{t|t-1}^{(k)}\}_{k=1}^{100}\right)\right].
\end{align}
因此，DistPred的輸出數、損失函數及區間定義固定，ANN與
Minusformer只負責提供不同的高頻特徵學習架構。

\subsection*{3.8.4 分位數回歸（Quantile Regression）}

Quantile方法直接估計$\tau\in\{0.025,0.5,0.975\}$三個條件分位數，
而不先產生大量預測members。對分位數水準$\tau$，令
$e_t=r_t-\widehat Q_{\tau,t}$，pinball loss為
\begin{equation}
 \rho_\tau(e_t)=
 \begin{cases}
  \tau e_t, & e_t\geq0,\\
  (\tau-1)e_t, & e_t<0.
 \end{cases}
\end{equation}
模型以三個分位數loss的平均值作為訓練目標。高分位數對低估給予
較大懲罰，低分位數則對高估給予較大懲罰，因此各輸出會分別學習
條件分布的中間位置與兩側尾端。

若三個分位數完全獨立輸出，可能發生下分位數高於中位數或上分位
數的quantile crossing。為避免此問題，令輸出層raw values為
$a_t,b_t,c_t$，本研究將其轉換為
\begin{align}
 \widehat Q_{0.5,t}&=a_t,\\
 \widehat Q_{0.025,t}&=a_t-\operatorname{Softplus}(b_t),\\
 \widehat Q_{0.975,t}&=a_t+\operatorname{Softplus}(c_t),
\end{align}
由於Softplus輸出非負值，可從模型結構上維持
$\widehat Q_{0.025,t}\leq\widehat Q_{0.5,t}
\leq\widehat Q_{0.975,t}$。0.5分位數作為高頻點預測，另外兩個
分位數作為高頻區間端點。使用ANN或Minusformer時，皆採相同的
分位數水準、pinball loss與monotone輸出，差別僅在骨幹網路。

\subsection*{3.8.5 Sort Bootstrap}

Sort Bootstrap不由單一ANN直接輸出預測分布，而是利用重複抽樣及重新訓練
形成模型間的預測分布。設高頻訓練序列長度為$n$，第$b$次先從
$\{0,\ldots,n-1\}$有放回抽取$n$個索引，再將索引由小至大排序：
\begin{equation}
 i_{(1)}^{(b)}\leq\cdots\leq i_{(n)}^{(b)}.
\end{equation}
依排序後索引形成
\begin{equation}
 r^{*(b)}=\left(r_{i_{(1)}^{(b)}},\ldots,
 r_{i_{(n)}^{(b)}}\right).
\end{equation}
排序可恢復被抽中觀察值的相對先後關係，但有放回抽樣仍會使部分
時間點重複、部分時間點未被抽中，也不代表完整保留原序列的實際
時間間距。

每組bootstrap序列使用長度30的滑動視窗建立輸入與下一期目標，
再依形成後的樣本順序以前80\%訓練、後20\%驗證。每組樣本均重新
初始化並以MSE訓練一個ANN，再對同一組testing輸入產生預測；此
程序重複$B=100$次。令第$b$個ANN的預測為
$\widehat r_{t|t-1}^{(b)}$，則
\begin{align}
 \widehat r_{t|t-1}
 &=\frac{1}{B}\sum_{b=1}^{B}\widehat r_{t|t-1}^{(b)},\\
 [\widehat L_t^{(H)},\widehat U_t^{(H)}]
 &=\left[Q_{0.025}
 \left(\{\widehat r_{t|t-1}^{(b)}\}_{b=1}^{B}\right),
 Q_{0.975}
 \left(\{\widehat r_{t|t-1}^{(b)}\}_{b=1}^{B}\right)\right].
\end{align}
因此，點預測取100組ANN預測的平均，區間則取其經驗分布的
2.5與97.5百分位數。Sort Bootstrap的預測分布來自100次重抽與
重訓，與DistPred由單一模型一次輸出100個members的方式不同。

\medskip
\noindent\textbf{高低頻結果整合。}

對B1至B5，令低頻ARIMA點預測及區間為
$\widehat\ell_{t|t-1}$與$[\widehat L_t^{(L)},\widehat U_t^{(L)}]$，
高頻結果為$\widehat r_{t|t-1}$與
$[\widehat L_t^{(H)},\widehat U_t^{(H)}]$。最終原始結果為
\begin{align}
 \widehat z_{t|t-1}
 &=\widehat\ell_{t|t-1}+\widehat r_{t|t-1},\\
 C_t^{Original}
 &=\left[\widehat L_t^{(L)}+\widehat L_t^{(H)},
          \widehat U_t^{(L)}+\widehat U_t^{(H)}\right].
\end{align}

\section*{3.9 B1至B5的Local-scale校準}

B1至B5的原始程式沒有可直接用於時間順序校準的樣本外殘差，
因此在模型外建立共同wrapper，不修改其模型、損失函數或原始區間。
具體流程如下：
\begin{enumerate}
 \item 將原training期依時間順序切為前80\% fitting prefix與後20\% calibration suffix；
 \item 僅以前80\%資料執行原方法，再預測後20\%，取得原始點預測與區間；
 \item 依第3.6.1節方式，利用calibration residual建立局部尺度$\widehat\sigma_i$；
 \item 對原始區間計算signed conformity score
 \begin{equation}
  u_i=\frac{\max\{\widehat L_i-z_i,z_i-\widehat U_i\}}
  {\widehat\sigma_i};
 \end{equation}
 \item 令
 \begin{equation}
  \gamma=\min\left\{1,
  \frac{\lceil(n_{cal}+1)(1-\alpha)\rceil}{n_{cal}}
  \right\},
 \end{equation}
 並取$q=Q_{\gamma}(\{u_i\})$；
 \item 以完整training資料重估原方法並產生testing原始區間，再校準為
 \begin{equation}
  C_t^{LS}=\left[
  \widehat L_t-\widehat\sigma_tq,
  \widehat U_t+\widehat\sigma_tq
  \right].
 \end{equation}
\end{enumerate}
當$q>0$時向外擴張，$q<0$時則可向內縮短。testing期間在發出當期
區間後才更新score pool，且pool長度維持與calibration suffix相同。
此wrapper與OOT-AEnbPI直接標準化OOT residual的方式不同，兩者不應混為
同一個校準流程。

\section*{3.10 評估指標}

點預測使用均方根誤差（RMSE）：
\begin{equation}
 RMSE=\sqrt{\frac{1}{H}\sum_{t=1}^{H}
 (z_t-\widehat z_{t|t-1})^2}.
\end{equation}
區間預測使用coverage、平均寬度、寬度中位數、下界失誤率、上界
失誤率與interval score。Coverage為
\begin{equation}
 Coverage=\frac{1}{H}\sum_{t=1}^{H}
 \mathbf 1(\widehat L_t\leq z_t\leq\widehat U_t).
\end{equation}
令$W_t=\widehat U_t-\widehat L_t$，平均寬度與中位數為
\begin{align}
 \overline W&=\frac{1}{H}\sum_{t=1}^{H}W_t,\\
 \widetilde W&=\operatorname{median}\{W_1,\ldots,W_H\}.
\end{align}
上下界失誤率為
\begin{align}
 Lower\ miss&=\frac{1}{H}\sum_{t=1}^{H}
 \mathbf 1(z_t<\widehat L_t),\\
 Upper\ miss&=\frac{1}{H}\sum_{t=1}^{H}
 \mathbf 1(z_t>\widehat U_t).
\end{align}
Interval score定義為
\begin{equation}
 IS_{\alpha}(l,u;z_t)=(u-l)
 +\frac{2}{\alpha}(l-z_t)\mathbf 1(z_t<l)
 +\frac{2}{\alpha}(z_t-u)\mathbf 1(z_t>u).
\end{equation}
RMSE、區間寬度與interval score愈小愈佳；coverage應接近95\%名目水準。
